% =====================================================================
% RESULTS & DISCUSSION — revised 2026-07-08. Structure follows your
% draft (\subsection{Cross-Resolution Comparison} + \subsection{Discussion}).
% Every number verified against exp/combined_resolution/{testset_table.md,
% resolution_table.md, inference_memory_table.md} and the per-epoch
% metrics in exp/bestcombo_resolutions_06-24_11-49 (SOURCE comments inline).
% Cross-refs used from the main file: tab:final_configs,
% eq:logpara_scaling, sec:eval_metrics, sec:tuning_protocol, sec:intro.
%
% STILL TO DO (nothing blocks the text except #1):
%  [ ] 1. REDO fig:mae_vs_resolution — image.png is a screenshot with
%         wrong values. Regenerate a log-log test-MAE-vs-resolution
%         figure from exp/combined_resolution/testset_table.md.
%  [ ] 2. Cite the NVIDIA Jetson Nano datasheet where marked and
%         double-check the specs against it.
%  [ ] 3. OPTIONAL: FLOPs per network x resolution (ptflops/thop, no
%         training) if you want FLOPs numbers instead of future work.
%  [ ] 4. Pick the final learning-curves file (this draft uses
%         mae_curves_by_resolution_thesis_2x3.png).
% =====================================================================

\subsection{Cross-Resolution Comparison}
\label{sec:results_resolution}

The tuned configurations (Table~\ref{tab:final_configs}) were
retrained from scratch at each of the six resolution variants, with
the density scale $k$ adjusted per resolution by
Eq.~\eqref{eq:logpara_scaling} and every other hyperparameter frozen
at its $386\times260$ optimum. Table~\ref{tab:resolution_test} reports
the results; Figure~\ref{fig:mae_vs_resolution} plots the trend and
Figure~\ref{fig:res_learning_curves} the underlying learning curves.

\begin{table}[t]
    \centering
    \caption{Counting error across input resolutions, in larvae per
    image (metrics as defined in Section~\ref{sec:eval_metrics}). Each
    cell's best-validation-MAE checkpoint is evaluated exactly once on
    the test split ($112$ images). The asterisk marks the tuning
    resolution; bold marks the best value per column. MobileCount's
    test differences across the three largest resolutions are within
    run-to-run noise (see text).}
    \label{tab:resolution_test}
    \small
    \begin{tabular}{l rrr rrr}
        \toprule
        &\multicolumn{3}{c}{CSRNet}&\multicolumn{3}{c}{MobileCount}\\
        \cmidrule(lr){2-4}\cmidrule(lr){5-7}
        Resolution & val MAE & test MAE & test MSE & val MAE & test MAE & test MSE \\
        \midrule
        $49\times33$     & $5.96$  & $5.46$  & $6.22$  & $2.68$  & $2.92$  & $3.96$ \\
        $97\times65$     & $1.03$  & $1.36$  & $2.40$  & $1.91$  & $2.00$  & $2.95$ \\
        $193\times130$   & $0.534$ & $0.885$ & $1.83$  & $0.861$ & $0.895$ & $1.25$ \\
        $386\times260^*$ & $0.288$ & $0.470$ & $0.973$ & $\mathbf{0.385}$ & $0.433$ & $0.653$ \\
        $772\times519$   & $0.193$ & $0.289$ & $0.683$ & $0.434$ & $0.471$ & $0.618$ \\
        $1544\times1038$ & $\mathbf{0.183}$ & $\mathbf{0.233}$ & $\mathbf{0.479}$ & $0.428$ & $\mathbf{0.361}$ & $\mathbf{0.521}$ \\
        \bottomrule
    \end{tabular}
    % SOURCE: exp/combined_resolution/testset_table.md — all values
    % verified 2026-07-08.
\end{table}

Both networks tolerate downscaling gracefully at first: from the
maximum down to $193\times130$ — a $64$-fold reduction in pixel count
— the test MAE stays below $0.9$ for both. Below that the degradation
is steep. At $49\times33$ the test MAE reaches $5.46$ (CSRNet) and
$2.92$ (MobileCount); expressed relative to the mean count of $47.9$
larvae per image, the best models miscount by ${\approx}0.5$--$1\%$,
whereas the smallest variant is off by $6$--$11\%$. Even the collapsed
variants remain a factor of $4$--$8$ better than the trivial
constant-prediction floor of ${\approx}23$ identified in
Section~\ref{sec:tuning_protocol}, but their relative error is an
order of magnitude above every other operating point. Notably, the
architecture ranking inverts at the bottom of the range: at
$49\times33$ MobileCount is nearly twice as accurate as CSRNet,
plausibly because CSRNet's $H/8$ output resolution reduces a
$49\times33$ input to a density map of roughly $6\times4$ cells.

\begin{figure}
    \centering
    \includegraphics[width=0.5\linewidth]{image.png}
    % TODO(#1): REDO — current file is a screenshot with wrong values.
    \caption{Test MAE against input resolution for both networks
    (log--log axes). The asterisked point marks the tuning resolution.}
    % TODO: re-check caption (axes/scale) once the figure is redone.
    \label{fig:mae_vs_resolution}
\end{figure}

\begin{figure*}
    \centering
    \includegraphics[width=\linewidth]{images/mae_curves_by_resolution_thesis_2x3.png}
    \caption{Validation MAE over training for the tuned configurations
    at all six resolution variants. Curves are truncated at epoch
    $1200$ for legibility; no run changes materially in the final
    $300$ epochs.}
    \label{fig:res_learning_curves}
\end{figure*}

\paragraph{CSRNet improves monotonically with resolution,}
reaching a test MAE of $0.233$ at $1544\times1038$ — the best result
in this thesis. The final doubling, however, buys only
$0.289\rightarrow0.233$ test MAE at roughly triple the training cost
(${\approx}10$\,h $\rightarrow$ ${\approx}29$\,h), placing the
accuracy--compute knee at $772\times519$.
% SOURCE: duration_sec 36,615 s / 105,157 s in
% bestcombo_resolutions_results.json (verified 2026-07-08).

\paragraph{MobileCount saturates at $386\times260$.}
MobileCount is non-monotonic: its validation optimum lies at the
tuning resolution ($0.385$), and the two larger variants are no better
($0.434$ and $0.428$); the test MAE fluctuates between $0.361$ and
$0.471$ across the three largest resolutions, differences within
plausible run-to-run variation (see the limitations in
Section~\ref{sec:discussion}). The learning curves
(Figure~\ref{fig:res_learning_curves}) reveal a marked difference in
training behaviour: at the two largest resolutions MobileCount reaches
its best checkpoint within the first ${\sim}250$ epochs and spends the
remainder of the budget at a validation MAE $3.7$--$3.8\times$ worse,
whereas CSRNet stays within ${\sim}1.3\times$ of its optimum to the
end. The interpretation of this behaviour is taken up in
Section~\ref{sec:discussion}.
% SOURCE: best epochs 105 (772x519) / 224 (1544x1038); ep>1200 median
% val MAE 1.60 vs best 0.43 for MobileCount; CSRNet 0.24 vs 0.19/0.18.

% ---------------------------------------------------------------------

\subsection{Discussion}
\label{sec:discussion}

\paragraph{MobileCount's saturation is capacity-limited overfitting.}
The learning curves suggest a mechanism for the saturation observed in
Section~\ref{sec:results_resolution}. At the two largest resolutions
MobileCount extracts what its $3.4$\,M parameters can represent from
the higher-resolution input within the first few hundred epochs; the
${\sim}1{,}300$ epochs that follow are spent memorising the $840$
training images rather than learning generalisable features, and its
validation MAE settles at $3.7$--$3.8\times$ its optimum.
Best-checkpoint selection shields the reported numbers from the
overfitting itself, but it cannot manufacture accuracy the network
never reaches. CSRNet — with $4.8\times$ the parameters and an
ImageNet-pretrained frontend — shows no such collapse: it remains
within ${\sim}1.3\times$ of its optimum throughout and keeps
converting additional resolution into accuracy.

\begin{figure}
    \centering
    \includegraphics[width=0.5\linewidth]{images/accuracy vs memory.png}
    \caption{Accuracy--memory frontier: test MAE against peak
    inference memory (batch-$1$, CUDA allocated peak) for all twelve
    network~$\times$~resolution cells. Labels give the input
    resolution.}
    % TODO: re-check caption against the final figure.
    \label{fig:acc_vs_mem}
\end{figure}

\paragraph{The accuracy--memory frontier crosses over at
${\sim}0.36$ MAE.} The shape of the frontier
(Figure~\ref{fig:acc_vs_mem}) follows from what dominates the
footprint. At small inputs it is the weights: the peaks sit nearly
flat at ${\approx}150$\,MB for CSRNet and ${\approx}27$\,MB for
MobileCount, tracking their parameter memory ($127.6$\,MB for
$16.3$\,M parameters vs.\ $26.3$\,MB for $3.4$\,M); from
$386\times260$ upwards, activations take over and the footprint grows
with pixel count. Across the frontier, MobileCount offers the better
accuracy-per-megabyte everywhere down to a test MAE of $0.361$
($217$\,MB at $1544\times1038$). At the $386\times260$ operating point
it strictly dominates CSRNet: better accuracy ($0.433$ vs.\ $0.470$
test MAE, ${\approx}0.9\%$ of the mean count) at roughly one-sixth the
memory ($39$\,MB vs.\ $229$\,MB). Below ${\sim}0.36$ MAE the frontier
belongs to CSRNet alone — $0.289$ at $524$\,MB and $0.233$ at
$1713$\,MB — and the step onto its branch is expensive: CSRNet's best
point costs ${\approx}44\times$ the memory of MobileCount's
$386\times260$ operating point (${\approx}8\times$ MobileCount at the
same resolution) and ${\approx}29\times$ the training time
(${\approx}29$\,h vs.\ ${\approx}1$\,h).
% SOURCE: exp/combined_resolution/inference_memory_table.md; weights
% 127.6/26.3 MB, params 16.26M/3.40M. 1713/39=44, 1713/217=7.9.

\paragraph{Implications for edge deployment.}
The intended deployment class makes the crossover concrete. An
embedded module such as the NVIDIA Jetson Nano provides a $128$-core
Maxwell GPU and $4$\,GB of LPDDR4 \emph{shared} between CPU and GPU,
of which the operating system and runtime consume a substantial
fraction.
% TODO(#2): cite the NVIDIA Jetson Nano datasheet and verify specs.
MobileCount's $39$\,MB peak at $386\times260$ is negligible on such a
device, whereas CSRNet's $524$\,MB at $772\times519$ is uncomfortable
and its $1713$\,MB at $1544\times1038$ is marginal at best. This
resolves the question posed in Section~\ref{sec:intro}: the ranking of
the two strategies is not fixed but conditional on the accuracy
requirement. If a counting error of ${\approx}1\%$ of the mean count
suffices, the lightweight model at $386\times260$ is the clear choice
for edge deployment; only when sub-$0.36$ MAE is required does the
heavy model at high resolution justify its footprint — and then likely
not on this device class at all.

\paragraph{Limitations.}
All results rest on a single random seed ($3035$). Three-seed
replicates of representative configurations, run during tuning, put
the across-seed validation spread at the tuning resolution at
$\pm0.007$ (CSRNet) to $\pm0.055$ (MobileCount) MAE; spreads at other
resolutions and on the test split were not measured, a deliberate
economy. The differences among MobileCount's three largest resolutions
are of this order, so only the qualitative shapes — CSRNet's monotone
improvement, MobileCount's saturation — should be considered robust.
% NOTE: the three-seed replicate numbers aren't stated anywhere in the
% main file anymore; consider adding a sentence to the tuning protocol.
% SOURCE: exp/optimizer_seed_sweep_* — CSRNet AdamW ±0.007,
% MobileCount ±0.055.
Hyperparameters were tuned at $386\times260$ only and reused at all
resolutions; a per-resolution tuning could shift individual numbers,
though at prohibitive cost. In particular, MobileCount's saturation is
confounded with the choice of tuning resolution — its hyperparameters
were selected at the very resolution where it peaks — although CSRNet,
tuned under the same handicap, improves at every resolution
regardless. Relatedly, the batch size was tuned conditionally on the
optimiser~$\times$~learning-rate optimum rather than as a joint grid,
so interactions along that axis cannot be fully excluded. The two
architectures also differ in more than size: initialisation
(ImageNet-pretrained frontend vs.\ training from scratch),
normalisation (none vs.\ batch normalisation), and design (a dilated
single-branch backend vs.\ an encoder--decoder), so their divergent
scaling behaviour cannot be attributed to capacity alone. The memory
figures are CUDA allocated peaks measured on a single GPU model; the
allocator's reserved pool is somewhat higher, and other hardware or
runtimes will shift the absolute values, though not the ranking
between the architectures. Finally, all experiments concern a single
dataset from one imaging setup, and the evaluation partition is drawn
from the same acquisition as the training data: the reported errors
characterise the curated setup, not production-farm imagery, and
transferability to other larval densities, camera geometries, or
species is untested.

\paragraph{Future work.}
Three extensions follow directly. First, the compute side of the
trade-off could be characterised alongside memory by measuring FLOPs
and, more importantly, actual latency per network and resolution.
% TODO(#3, optional): ptflops/thop pass if you want FLOPs numbers now.
Second, the deployment argument should be closed by running the
selected configuration on the target hardware itself — a Jetson-class
module — measuring end-to-end latency and peak memory, and assessing
quantisation. Third, the setup should be validated under realistic
domain shift: images from operating farms, where robustness-oriented
augmentation beyond the patch-based scheme used here may earn its
keep.
